\clearpage
\phantomsection
\addcontentsline{toc}{section}{References}
\begingroup
\small
\setlength{\parskip}{0pt}
\begin{thebibliography}{99}
\setlength{\itemsep}{4pt}
\raggedright

\bibitem{RepoSnapshot}
V. Reshetnikov, \emph{ProveIt: Analysis/Polylogarithms},
repository snapshot
\texttt{3a6d80ed618d839915deb0d19ab687f45e81d995},
8 October 2026 UTC.
\url{https://github.com/VladimirReshetnikov/ProveIt/tree/3a6d80ed618d839915deb0d19ab687f45e81d995/Analysis/Polylogarithms/docs}.
The intake README is in the parent directory; source paths and Git
blob identifiers are listed in the accompanying \path{PROVENANCE.json}.

\bibitem{DraftTower}
ProveIt draft corpus, \emph{The Parameter-Derivative Tower of the
Generalized Stieltjes Constants: a Uniform Harmonic-Number Theorem}, source
\path{docs/articles/stieltjes-parameter-derivative-tower.tex},
in the snapshot \cite{RepoSnapshot}.

\bibitem{DraftLadder}
ProveIt draft corpus, \emph{Stieltjes Antiderivatives, Hurwitz Zeta Jets,
and Polygamma Functions of Negative Order}, source
\path{docs/articles/stieltjes-antiderivative-ladder.tex},
in the snapshot \cite{RepoSnapshot}.

\bibitem{DLMF25}
NIST, \emph{Digital Library of Mathematical Functions},
Chapter 25: Zeta and Related Functions, especially
\S25.11 (Hurwitz zeta) and \S25.15 (Dirichlet $L$-functions).
\url{https://dlmf.nist.gov/25.11};
\url{https://dlmf.nist.gov/25.15}.
Accessed 8 October 2026.

\bibitem{DLMF5}
NIST, \emph{Digital Library of Mathematical Functions},
Chapter 5: Gamma Function, especially \S5.4 (special values),
\S5.5 (functional relations), and \S5.8 (infinite products).
\url{https://dlmf.nist.gov/5.4};
\url{https://dlmf.nist.gov/5.5};
\url{https://dlmf.nist.gov/5.8}.
Accessed 8 October 2026.

\bibitem{Coffey2009}
M. W. Coffey, “Series representations for the Stieltjes constants,”
\emph{Rocky Mountain Journal of Mathematics} \textbf{44} (2014),
443--477; preprint arXiv:0905.1111 (2009).
\url{https://arxiv.org/abs/0905.1111}.
Proposition 5 treats the unique maximum of $\gamma_1(a)$.

\bibitem{Kubert1979}
D. S. Kubert, “The universal ordinary distribution,”
\emph{Bulletin de la Soci\'et\'e Math\'ematique de France}
\textbf{107} (1979), 179--202.
\url{https://www.numdam.org/item/BSMF_1979__107__179_0.pdf}.

\bibitem{gau-olaikhan-2022}
A. Olaikhan, answers on Mathematics Stack Exchange:
“How to write a limit in terms of finite summation,”
3 March 2022, edited 5 March 2022,
\url{https://math.stackexchange.com/questions/4139313/how-to-write-a-limit-in-terms-of-finite-summation/4395416};
and the all-weight generalization posted on 5 March 2022 to
“Evaluate $\sum_{n=1}^{\infty}(-1)^{n-1}H_n/(2n+1)^3$,”
\url{https://math.stackexchange.com/questions/3262663/evaluate-sum-n-1-infty-frac-1n-1h-n2n13}.

\bibitem{audit-glanois-2016}
C. Glanois, “Motivic unipotent fundamental groupoid of
$\mathbb G_m\setminus\mu_N$ for $N=2,3,4,6,8$ and Galois descents,”
\emph{Journal of Number Theory} \textbf{160} (2016), 334--384.
\href{https://doi.org/10.1016/j.jnt.2015.08.003}{doi:10.1016/j.jnt.2015.08.003};
\url{https://arxiv.org/abs/1411.4947}, especially \S3.1 of version 2.

\bibitem{audit-pari-polylog}
The PARI Group, \emph{PARI/GP documentation: Transcendental functions},
entry \texttt{polylog}, flag $3$ (Zagier's modified polylogarithm).
\url{https://pari.math.u-bordeaux.fr/dochtml/html-stable/Transcendental_functions.html}.
Accessed 8 October 2026.

\bibitem{EspinosaMoll2002b}
O. Espinosa and V. H. Moll, “On Some Integrals Involving the
Hurwitz Zeta Function: Part 2,” \emph{Ramanujan Journal}
\textbf{6} (2002), 449--468.
\url{https://www.math.tulane.edu/~vhm/papers_html/hurwitz2.pdf}.

\bibitem{RadchenkoZagier2023}
D. Radchenko and D. Zagier, “Arithmetic properties of the Herglotz
function,” \emph{Journal f\"ur die reine und angewandte Mathematik}
\textbf{797} (2023), 229--253.
\url{https://arxiv.org/abs/2012.15805}.

\bibitem{MilneCFT}
J. S. Milne, \emph{Class Field Theory}, version 4.03,
6 August 2020.
\url{https://www.jmilne.org/math/CourseNotes/CFT.pdf}.

\bibitem{MilneANT}
J. S. Milne, \emph{Algebraic Number Theory}, version 3.08,
19 July 2020, especially Theorem 4.3 and
Proposition 7.55 with Remark 7.56.
\url{https://www.jmilne.org/math/CourseNotes/ANT.pdf}.

\bibitem{RivoalRoques2014}
T. Rivoal and J. Roques, \emph{Hadamard products of algebraic functions},
author preprint, 29 April 2014.
\url{https://rivoal.perso.math.cnrs.fr/articles/algcoef.pdf}.
See the discussion of gamma relations and Rohrlich's conjecture.

\end{thebibliography}
\endgroup
\end{document}
